% P11 paper module: R15 exact gauge-intertwining criterion and square-root information gate.

\subsection{Gauge intertwining and the finite-jet square-root information gate}
\label{subsec:o3n-gauge-intertwining}

Keep $0<R<S<T_0<U$ fixed on the odd sectors and retain the notation
\[
 W:=W_{R,S,-}^{[T_0]},\qquad
 W_U:=W_{R,S,-}^{[U]},
\]
\[
 X_R=U_RA_R^{1/2},\qquad X_S=U_SA_S^{1/2},
\qquad
 Q=A_S^{1/2}WA_R^{-1/2},
\]
from Remarks~\ref{rem:polar-gauge} and Subsection~\ref{subsec:o3m-polar-gauge}.  Recall
\[
 W_U=U_SQU_R^*,
 \qquad
 \Gamma_U=W^*U_SWU_R^*.
\]

\begin{proposition}[Exact gauge-intertwining criterion]\label{prop:o3n-gauge-criterion}
Put
\[
 V_U:=U_SWU_R^*.
\]
Then $V_U$ and $W$ are isometries and
\[
 \Gamma_U=W^*V_U.
\]
For every source vector $f$,
\[
 \boxed{
 \|(V_U-W)f\|^2
 =2\|f\|^2-2\operatorname{Re}\langle\Gamma_Uf,f\rangle.
 }
\tag{O3N.1}\label{eq:o3n-gauge-rayleigh}
\]
Consequently
\[
 \boxed{
 \Gamma_U\xrightarrow[s]{}I
 \quad\Longleftrightarrow\quad
 V_U\xrightarrow[s]{}W.
 }
\tag{O3N.2}\label{eq:o3n-gauge-strong}
\]
In operator norm,
\[
 \boxed{
 \|\Gamma_U-I\|
 \le \|V_U-W\|
 \le \sqrt{2\|\Gamma_U-I\|},
 }
\tag{O3N.3}\label{eq:o3n-gauge-norm}
\]
so norm convergence of $\Gamma_U$ to $I$ is equivalent to norm convergence of the
transported gauge isometry $V_U$ to $W$.  Equivalently in operator norm,
\[
 \boxed{
 \Gamma_U\to I
 \quad\Longleftrightarrow\quad
 \|U_SW-WU_R\|\to0.
 }
\tag{O3N.4}\label{eq:o3n-gauge-intertwining}
\]
\end{proposition}

\begin{proof}
Since $U_R,U_S$ are unitary and $W$ is an isometry, so is $V_U$.  Therefore
\[
 (V_U-W)^*(V_U-W)=2I-W^*V_U-V_U^*W
 =2I-\Gamma_U-\Gamma_U^*,
\]
which gives~\eqref{eq:o3n-gauge-rayleigh}.  If $\Gamma_Uf\to f$ for a fixed $f$, the
right-hand side tends to zero, proving $V_Uf\to Wf$.  The converse follows from
$\Gamma_U-I=W^*(V_U-W)$.  This proves~\eqref{eq:o3n-gauge-strong}.

The same identity gives
\[
 \|V_U-W\|^2
 \le 2\|I-\Gamma_U\|,
\]
while $\|\Gamma_U-I\|\le\|V_U-W\|$, proving~\eqref{eq:o3n-gauge-norm}.  Finally
\[
 V_U-W=(U_SW-WU_R)U_R^*,
\]
so the two operator norms are equal.
\end{proof}

The preceding proposition isolates the genuine gauge question.  The unnormalised
square-root intertwining suggested by the metric products contains both the modulus and
gauge defects and is stronger than necessary.

\begin{proposition}[Exact square-root coherence decomposition]\label{prop:o3n-X-decomp}
One has
\[
 \boxed{
 X_SW-WX_R
 =\Bigl[U_S(Q-W)+(U_SW-WU_R)\Bigr]A_R^{1/2}.
 }
\tag{O3N.5}\label{eq:o3n-X-decomp}
\]
Equivalently, after the canonical right normalisation,
\[
 \boxed{
 \mathcal E_U
 :=(X_SW-WX_R)A_R^{-1/2}
 =U_SQ-WU_R.
 }
\tag{O3N.6}\label{eq:o3n-X-normalized}
\]
Moreover
\[
 \boxed{
 \mathcal E_UU_R^*=W_U-W,
 \qquad
 \|\mathcal E_U\|=\|W_U-W\|.
 }
\tag{O3N.7}\label{eq:o3n-X-transport}
\]
Thus the correctly normalised $X_S$--$X_R$ intertwining problem is exactly the original
norm terminal-transport problem; it is not an independent simplification.
\end{proposition}

\begin{proof}
From $A_S^{1/2}W=QA_R^{1/2}$ and the polar decompositions,
\[
 X_SW=U_SA_S^{1/2}W=U_SQA_R^{1/2},
 \qquad
 WX_R=WU_RA_R^{1/2}.
\]
Subtracting gives~\eqref{eq:o3n-X-decomp} and~\eqref{eq:o3n-X-normalized}.  Finally
$W_U=U_SQU_R^*$ gives~\eqref{eq:o3n-X-transport}.
\end{proof}

There is also an exact local interpretation of each individual polar gauge.

\begin{proposition}[Individual polar gauge equals metric noncommutativity]
\label{prop:o3n-gauge-commutator}
Let $X\in\{R,S\}$ and abbreviate
\[
 B:=G_{X,T_0}^-,\qquad C:=G_{X,U}^-.
\]
For the polar factor $U_X$ of
\[
 C^{1/2}B^{-1/2}=U_X
 \bigl(B^{-1/2}CB^{-1/2}\bigr)^{1/2},
\]
one has
\[
 \boxed{
 U_X=I
 \quad\Longleftrightarrow\quad
 [B,C]=0.
 }
\tag{O3N.8}\label{eq:o3n-gauge-commutator}
\]
Hence the word ``gauge'' records genuine noncommutativity between the fixed base metric
and the future metric; it is not a free phase choice.
\end{proposition}

\begin{proof}
If $B$ and $C$ commute, then $C^{1/2}B^{-1/2}$ is positive, so its polar unitary is
$I$.  Conversely, if $U_X=I$, then $C^{1/2}B^{-1/2}$ equals its positive polar modulus
and is therefore selfadjoint.  Hence
\[
 C^{1/2}B^{-1/2}=B^{-1/2}C^{1/2},
\]
so the positive square roots commute, equivalently $B$ and $C$ commute.
\end{proof}

The mixed-jet theorem determines all fixed-pair leading matrix elements, but its
leading matrix is asymptotically rank one.  The next elementary model shows why the
unquantified $o(1)$ remainders are insufficient for inverse-square-root control.

\begin{proposition}[Rank-one leading data do not determine inverse square roots]
\label{prop:o3n-rankone-insufficiency}
Let $\varepsilon\downarrow0$ and, for any fixed $a>2$, define
\[
 v_\varepsilon:=\binom{1}{\varepsilon},
 \qquad
 M_\varepsilon^{(a)}
 :=v_\varepsilon v_\varepsilon^*
   +\varepsilon^a e_2e_2^*
 =\begin{pmatrix}
 1&\varepsilon\\
 \varepsilon&\varepsilon^2+\varepsilon^a
 \end{pmatrix}>0.
\]
For every $a>2$ these matrices have exactly the same entrywise leading asymptotics
\[
 (M_\varepsilon^{(a)})_{11}=1,
 \qquad
 (M_\varepsilon^{(a)})_{12}=\varepsilon,
 \qquad
 (M_\varepsilon^{(a)})_{22}=\varepsilon^2(1+o(1)).
\tag{O3N.9}\label{eq:o3n-same-leading}
\]
However
\[
 \boxed{
 \lambda_{\min}(M_\varepsilon^{(a)})\sim\varepsilon^a,
 \qquad
 \|(M_\varepsilon^{(a)})^{-1/2}\|
 \sim\varepsilon^{-a/2}.
 }
\tag{O3N.10}\label{eq:o3n-inverse-root-scale}
\]
Thus identical rank-one fixed-pair leading data, even with relative $o(1)$ in every
entry, do not determine the polynomial scale of the inverse square root.
\end{proposition}

\begin{proof}
The determinant and trace are
\[
 \det M_\varepsilon^{(a)}=\varepsilon^a,
 \qquad
 \operatorname{tr}M_\varepsilon^{(a)}
 =1+\varepsilon^2+\varepsilon^a=1+o(1).
\]
Hence the larger eigenvalue tends to $1$, while the smaller eigenvalue, equal to the
determinant divided by the larger one, is asymptotic to $\varepsilon^a$.  The
inverse-square-root norm is the reciprocal square root of the smallest eigenvalue.
\end{proof}

The TC1 decomposition identifies the corresponding near-null direction inside the
actual fixed two-vector jet block.

\begin{proposition}[Exact TC1 near-null direction]\label{prop:o3n-nearnull}
Fix smooth odd vectors $f_0,f_m$ at one source level, with first nonzero integral-jet
orders $0$ and $m>0$, respectively.  For all sufficiently large $U$ one has
$\ell_U(f_0)\ne0$.  Define the $U$-dependent vector
\[
 \boxed{
 z_U:=f_m-\frac{\ell_U(f_m)}{\ell_U(f_0)}f_0.
 }
\tag{O3N.11}\label{eq:o3n-zU}
\]
Then the constant-mode rank-one form vanishes exactly on $z_U$,
\[
 \boxed{
 \rho_U(z_U,z_U)=0,
 \qquad
 \sigma_U(Jz_U,Jz_U)=D_U(z_U,z_U).
 }
\tag{O3N.12}\label{eq:o3n-nearnull-exact}
\]
Moreover
\[
 \boxed{
 D_U(z_U,z_U)
 =o\!\left(\frac{e^U}{U^{2m+2}}\right).
 }
\tag{O3N.13}\label{eq:o3n-nearnull-littleo}
\]
The present theory supplies no matching positive lower asymptotic for this quantity.
\end{proposition}

\begin{proof}
By~\eqref{eq:mj-ell},
\[
 \frac{\ell_U(f_m)}{\ell_U(f_0)}=O(U^{-m}),
\]
and $\ell_U(f_0)\ne0$ for all sufficiently large $U$.  The definition
\eqref{eq:o3n-zU} gives $\ell_U(z_U)=0$, hence
\[
 \rho_U(z_U,z_U)=\frac{|\ell_U(z_U)|^2}{d_U}=0.
\]
Since $\sigma_U=\rho_U+D_U$ on the fixed span of $f_0,f_m$, this proves
\eqref{eq:o3n-nearnull-exact}.

Expand $D_U(z_U,z_U)$ sesquilinearly in the fixed pair $f_0,f_m$.  Equations
\eqref{eq:mj-Ddiag}--\eqref{eq:mj-Dmixed}, together with
$\ell_U(f_m)/\ell_U(f_0)=O(U^{-m})$, show that each resulting term is
$o(e^U/U^{2m+2})$.  This proves~\eqref{eq:o3n-nearnull-littleo}.
\end{proof}

\begin{remark}[R15 information firewall]\label{rem:o3n-firewall}
On a two-vector jet-adapted block with first jet orders $0$ and $m>0$, the mixed-jet
asymptotic, after removal of its common scalar factor and a fixed diagonal
normalisation by the nonzero leading jet coefficients, has the rank-one scale
\[
 \begin{pmatrix}
 1&U^{-m}\\
 U^{-m}&U^{-2m}
 \end{pmatrix}
\]
up to entrywise relative $o(1)$ errors.  Proposition~\ref{prop:o3n-rankone-insufficiency}
shows that such information does not determine the smallest eigenvalue or the inverse
square root even to polynomial order.  Proposition~\ref{prop:o3n-nearnull} identifies
the exact P11 two-jet direction on which the leading rank-one term cancels: the missing
scale is the positive TC1 Gram remainder $D_U(z_U,z_U)$ (together with the fixed Gamma
floor in the full metric), not another copy of the already known leading boundary jet.

Therefore the present fixed-pair asymptotics cannot by themselves decide $U_R,U_S$,
$\Gamma_U$, or the true cross-terminal kernel.  This is an information-theoretic
obstruction to the current proof package, not a counterexample to the actual P11 metric
family.

Accordingly the genuine next step is exactly Open Problem~\ref{open:finite-jet-sqrt}:
one needs a quantitative finite-dimensional jet expansion deep enough to resolve the
near-null Gram directions and then compare the resulting square roots between the
$R$- and $S$-levels.  At the two-jet level, the sharpest immediate target is a
quantitative asymptotic or two-sided scale for $D_U(z_U,z_U)$.  No conclusion about
$\Gamma_U\to I$, $K_{R,S}^{T_0,U}\to I$, strong terminal transport, Object X, or RH is
drawn here.
\end{remark}
